\begin{proof}
\textbf{1. Reduction to the case \(a,b>0\).}  
If \(a=0\) or \(b=0\) then the left–hand side is \(0\) and the right–hand side is
non‑negative, so the inequality holds. Hence we may assume \(a>0,\;b>0\).

\medskip
\textbf{2. Introduce the conjugate exponent.}  
Let \(p>1\) and define \(q\) by  
\[
\frac{1}{p}+\frac{1}{q}=1\qquad\Longrightarrow\qquad 
q=\frac{p}{p-1}>1 .
\]

\medskip
\textbf{3. Convexity of the exponential function.}  
The function \(\varphi(t)=e^{t}\) is convex on \(\mathbb{R}\).  
By Jensen’s inequality, for any \(\lambda\in[0,1]\) and any real numbers \(x,y\),

\begin{align}
e^{\lambda x+(1-\lambda)y}\le \lambda e^{x}+(1-\lambda)e^{y}.
\tag{1}
\end{align}

\medskip
\textbf{4. Choice of the convex combination.}  
Set
\[
\lambda=\frac{1}{p},\qquad 1-\lambda=\frac{1}{q},\qquad
x=p\ln a,\qquad y=q\ln b .
\]
Then
\begin{align}
\lambda x+(1-\lambda)y
&=\frac{1}{p}\,p\ln a+\frac{1}{q}\,q\ln b \\
&=\ln a+\ln b \\
&=\ln(ab).
\tag{2}
\end{align}

\medskip
\textbf{5. Apply Jensen’s inequality.}  
Substituting the choices of \(\lambda,x,y\) and using \((2)\) in \((1)\) gives

\begin{align}
e^{\ln(ab)}\le \frac{1}{p}e^{p\ln a}+\frac{1}{q}e^{q\ln b}.
\end{align}
Since \(e^{\ln(ab)}=ab\) and \(e^{p\ln a}=a^{p},\;e^{q\ln b}=b^{q}\), we obtain

\begin{align}
ab\le \frac{a^{p}}{p}+\frac{b^{q}}{q}.
\tag{3}
\end{align}

\medskip
\textbf{6. Extension to the non‑negative case.}  
Both sides of \((3)\) are continuous in \((a,b)\) on \([0,\infty)^{2}\); therefore the
inequality, proved for all \(a,b>0\), remains valid when either \(a=0\) or \(b=0\).

\medskip
\textbf{7. Equality condition.}  
Equality in Jensen’s inequality \((1)\) occurs exactly when the two arguments are
equal, i.e. when \(x=y\).  With the chosen \(x,y\) this means

\[
p\ln a = q\ln b \;\Longleftrightarrow\; a^{p}=b^{q}.
\]

Hence equality in Young’s inequality holds iff \(a^{p}=b^{q}\).

\end{proof}